% Conclusion for the short Boudourides submission only (SipArxiv.tex).
% The companion paper keeps conc.tex, which cites the SIS persistence,
% extinction and higher-order cusp results that this version omits.
\section{Conclusions and remaining problems}\label{s:concShort}

Siphon sets cut out invariant siphon faces. First-order terms on such a face
assemble its transverse invasion matrix, while terms at least quadratic in the
absent variables leave that matrix unchanged and can still create bistability
% p:KBsiph, p:KBresponsezeros and p:KBorder are three \label aliases on the one
% proposition in mod.tex, so cite only one of them or the number repeats.
(Proposition~\ref{p:KBsiph}). The transverse block is Metzler, so its spectral
abscissa determines linear invasion or decay when nonzero; the critical case
requires nonlinear analysis.

For threshold networks, Theorem~\ref{t:KBsharp} shows that the threshold $\bar R_0=1/L$ is
sharp over the response class with slope bound $L$.
Theorem~\ref{t:siphMulti} instead fixes the siphon lattice and every transverse
matrix, using responses flat at onset. In both, the response depends on the
prescribed count; neither gives arbitrary multistability for a fixed response or
a fixed polynomial degree.

Odds-polynomial responses join these directions. Their zero set fixes the siphon
lattice, their onset slope enters the transverse matrix directly, and their higher
coefficients enter it only through the resident state (Proposition~\ref{p:Odd}).
Theorem~\ref{t:Net} carries the quartic cusp of Proposition~\ref{p:Qua} into a
network over a transversely stable siphon face; singleton siphons force
multiplicative coupling. Overlapping minimal siphons allow additive
cross-couplings, and Theorem~\ref{t:KBoverlapCusp} combines them with nonzero
residents, distinct explicit invasion numbers and an interior cusp. This cusp has
a positive cubic coefficient: it is unstable, the middle equilibrium of its
unfolding is stable, the stable equilibria at the cusp parameter lie elsewhere,
and Corollary~\ref{c:Nin} adds one stable equilibrium to the eight.

Three problems remain. Local bifurcation analysis does not determine which
attractor is reached after a fold. Theorem~\ref{t:Net} gives no explicit coupling
window along its continued family, where $a_1$ moves with the couplings. Finally,
it remains open how the cusp branches intersect the invasion-threshold
hypersurfaces $R_{\Sigma}=1$ associated with overlapping minimal siphons.
Theorem~\ref{t:KBoverlapCusp} proves coexistence and supplies the thresholds, but
does not establish such intersections; the same question arises for asymmetric
directed coupling, with $\gamma$ and $\varepsilon$ independent rather than tied.
